%%=============================================================================
%% Samenvatting
%%=============================================================================

% TODO: De "abstract" of samenvatting is een kernachtige (~ 1 blz. voor een
% thesis) synthese van het document.
%
% Een goede abstract biedt een kernachtig antwoord op volgende vragen:
%
% 1. Waarover gaat de bachelorproef?
% 2. Waarom heb je er over geschreven?
% 3. Hoe heb je het onderzoek uitgevoerd?
% 4. Wat waren de resultaten? Wat blijkt uit je onderzoek?
% 5. Wat betekenen je resultaten? Wat is de relevantie voor het werkveld?
%
% Daarom bestaat een abstract uit volgende componenten:
%
% - inleiding + kaderen thema
% - probleemstelling
% - (centrale) onderzoeksvraag
% - onderzoeksdoelstelling
% - methodologie
% - resultaten (beperk tot de belangrijkste, relevant voor de onderzoeksvraag)
% - conclusies, aanbevelingen, beperkingen
%
% LET OP! Een samenvatting is GEEN voorwoord!

%%---------- Nederlandse samenvatting -----------------------------------------
%
% TODO: Als je je bachelorproef in het Engels schrijft, moet je eerst een
% Nederlandse samenvatting invoegen. Haal daarvoor onderstaande code uit
% commentaar.
% Wie zijn bachelorproef in het Nederlands schrijft, kan dit negeren, de inhoud
% wordt niet in het document ingevoegd.

\IfLanguageName{english}{%
\selectlanguage{dutch}
\chapter*{Samenvatting}
\lipsum[1-4]
\selectlanguage{english}
}{}

%%---------- Samenvatting -----------------------------------------------------
% De samenvatting in de hoofdtaal van het document

\chapter*{\IfLanguageName{dutch}{Samenvatting}{Abstract}}

De mainframe-sector, en in het bijzonder z/OS-omgevingen, kampt met een toenemend tekort aan gekwalificeerde professionals. Dit tekort wordt versterkt door de vergrijzing binnen de sector en de beperkte aansluiting tussen het onderwijs en de noden van de industrie. Voor rollen zoals z/OS System Administrators en System Programmers ontbreekt bovendien vaak een duidelijk en praktijkgericht competentieframework, terwijl deze functies cruciaal zijn voor het beheer en de continuïteit van bedrijfskritische systemen.

Deze bachelorproef vertrekt vanuit de onderzoeksvraag: “Welke competenties zijn vereist voor z/OS System Administrators en System Programmers, en hoe kunnen deze op een gestructureerde en praktijkgerichte manier worden samengebracht in een competentieframework dat toepasbaar is binnen een onderwijscontext?” Het doel van dit onderzoek is de ontwikkeling van een dergelijk framework dat aansluit bij de realiteit van het werkveld.

De gehanteerde methodologie combineert een literatuurstudie rond z/OS, bestaande competentiemodellen en relevante technologieën, met kwalitatief onderzoek in de vorm van semigestructureerde interviews met ervaren professionals uit de sector. Deze professionals waren werkzaam bij organisaties zoals ABN AMRO, KBC, Isbank, Planet NL en Broadcom. De interviews boden waardevolle inzichten in de vereiste competenties, de moeilijkheden bij instap in het domein en de evoluties binnen de sector.

Een belangrijk inzicht uit het onderzoek is dat de afbakening tussen System Administrator en System Programmer in de praktijk niet eenduidig is en sterk kan variëren per organisatie. Dit leidde tot een verbreding van de scope van het onderzoek, waarbij beide rollen gecombineerd werden binnen één geïntegreerd competentieframework.

Het resultaat van deze bachelorproef is een gestructureerd en praktijkgericht competentieframework dat zowel technische als niet-technische competenties omvat. Technische domeinen bevatten onder meer kennis van z/OS-concepten, JCL, REXX, SDSF, JES en beveiliging (RACF), terwijl ook aandacht wordt besteed aan moderniseringstechnologieën zoals Liberty en WebSphere. Daarnaast worden professionele vaardigheden zoals communicatie, probleemoplossend denken en samenwerking geïntegreerd.

De conclusie van dit onderzoek is dat een dergelijk geïntegreerd framework een waardevolle bijdrage kan leveren aan het verbeteren van opleidingen rond mainframe-technologie en het verkleinen van de kloof tussen onderwijs en industrie. Aanbevelingen omvatten de integratie van het framework in curricula, verdere validatie met industriële partners en uitbreiding naar verschillende ervaringsniveaus.
